\chaves{Engenharia de Software; Fábrica de Software; Autenticação e Segurança; Qualidade de Software; Desenvolvimento Full-Stack.}

\begin{resumo} 

\textbf{Contexto:}
A formação acadêmica em Engenharia de Software requer vivência prática para a consolidação do aprendizado teórico. Atuar em ambientes que simulam a indústria, como uma Fábrica de Software, proporciona uma transição realista, exigindo a adaptação a arquiteturas existentes, a garantia contínua da qualidade do produto e a adoção de protocolos de segurança modernos.

\textbf{Objetivo:}
Este trabalho tem como objetivo descrever e analisar a experiência de evolução e refatoração full-stack do Sistema Integrado de Processos Seletivos (SIPROS) da UFG, destacando a implementação de mecanismos robustos de segurança, melhorias na interface de usuário e a adoção de práticas de garantia da qualidade.

\textbf{Método:}
O trabalho foi desenvolvido sob a metodologia de pesquisa-ação no contexto da disciplina de Fábrica de Software no semestre 2026/1. Utilizaram-se métodos ágeis para o planejamento e desenvolvimento iterativo. A execução envolveu processos rigorosos de Verificação e Validação (V\&V) entre Casos de Uso e protótipos, além da aplicação sistemática de revisões por pares (\textit{Code Review}) e testes automatizados.

\textbf{Resultados:}
A intervenção resultou na integração bem-sucedida da autenticação do Google via Keycloak, adotando a extensão PKCE (\textit{Proof Key for Code Exchange}) para mitigar vulnerabilidades estruturais no \textit{front-end} em Angular. Simultaneamente, as práticas de V\&V permitiram a correção de dezenas de inconsistências de arquitetura e usabilidade, enquanto a estruturação formal da suíte de testes de unidade garantiu a integridade do \textit{back-end} contra regressões.

\textbf{Conclusões:}
Os resultados evidenciam que a implementação de protocolos de segurança em Aplicações de Página Única (SPAs) exige um rigor arquitetural que afeta todo o fluxo do sistema. A experiência atestou que o emprego constante de práticas de qualidade (testes e V\&V) é indispensável em projetos colaborativos. Por fim, a atuação na Fábrica de Software proporcionou um amadurecimento técnico e comportamental inestimável, aproximando genuinamente a vivência acadêmica da realidade do mercado de tecnologia.

\end{resumo}
